%% kapa: exam=1 round=35 year=2024 period=1 form=A subject=civil_law subject_ko=민법 seq=32 number=32
%% source: Q-Net 감정평가사 기출문제 (cst00309) artlSeq=5213229 file=raw/1008/5213229/1교시 문제지 원본.pdf
%% conversion: kordoc 4.24.51 -> tools/gen1.py (auto); review=proofread-sonnet; flags=none
\begin{question}{32}
X토지를 3분의 1씩 공유하는 甲, 乙, 丙의 공유물분할에 관한 설명으로 옳지 않은 것은? (다툼이 있으면 판례에 따름)
\begin{choices}
  \item 甲은 乙과 丙의 동의를 얻지 않고서 공유물의 분할을 청구할 수 있다.
  \item 甲, 乙, 丙이 3년간 공유물을 분할하지 않기로 합의한 것은 유효하다.
  \item 공유물분할의 소에서 법원은 X를 甲의 단독소유로 하고 乙과 丙에게 지분에 대한 합리적인 가액을 지급하도록 할 수 있다.
  \item 甲의 지분 위에 설정된 근저당권은 공유물분할이 되어도 특단의 합의가 없는 한 X 전부에 관하여 종전의 지분대로 존속한다.
  \item 甲, 乙, 丙 사이에 공유물분할에 관한 협의가 성립하였으나 분할협의에 따른 지분이전등기에 협조하지 않으면 공유물분할의 소를 제기할 수 있다.
\end{choices}
\end{question}
\answer{5}
